\begin{enumerate}
\item[\textbf{Q1.}]\textbf{Latency-vs-fidelity crossover on real dynamic-circuit hardware.}
Griffiths-Niu's Fig.~2 assumes ideal projective measurements with zero
delay before the classical signal reaches the next box. On real hardware,
mid-circuit measurement latency is ${\sim}1\,\mu$s and the remaining qubits
decohere during that window. What is the threshold ratio
$T_2 / (t_{\text{meas}} + t_{\text{classical\_routing}})$ at which the
semiclassical QFT begins to \emph{lose} fidelity relative to the standard
QFT, when we insert per-qubit depolarizing/$T_2$ dephasing during the
classical-feed-forward window? Our Aer replication used an ideal
zero-latency \texttt{if\_test} block, which trivially matches the standard
QFT --- the paper predates real dynamic-circuit hardware and so cannot
answer this.

\item[\textbf{Q2.}]\textbf{Shot-noise scaling of TVD(full$|$semi) with number of Fourier peaks.}
In our periodic-input tests ($n=4$, $p=8$) the semiclassical result showed
the largest full-vs-semi TVD ($0.011$) of any case --- roughly $3{\times}$
larger than the $p=2$ case at the same shot budget. Is this genuinely
extra shot-noise correlation between measurement outcomes and later
phases (in which case the semiclassical variant has \emph{strictly higher}
shot-noise variance than the standard variant), or a transpiler artifact
of the unrolled \texttt{if\_test} blocks?  Formalize: prove or disprove
that the joint distributions of the $n$ classical outcomes are
\emph{equal} as measures (the paper proves only marginals).

\item[\textbf{Q3.}]\textbf{Partially-semiclassical QFT for non-terminal QFT sub-blocks.}
Griffiths-Niu eliminate every two-qubit gate only because the QFT is
followed immediately by full measurement. In algorithms that use the
post-QFT state coherently (amplitude estimation without terminal
measurement, QSVT sub-blocks), the substitution does not apply as
stated. Is there a hybrid variant that measures only the first
$m<n$ qubits semiclassically --- converting $m(n{-}m)$ two-qubit gates
into classically-conditioned single-qubit gates, leaving the remaining
coherent register live?  What is the Pareto frontier of (2q-gate count,
state fidelity to ideal QFT on the $n{-}m$ remaining qubits) as $m$
varies?

\item[\textbf{Q4.}]\textbf{A channel-identity proof of equivalence, without consistent histories.}
The paper's argument uses the consistent-histories formalism.  Can the
equivalence be written directly as a CPTP-channel identity of the form
$(\text{measure then classical-controlled }P) = (\text{CP then measure})$
on the two-qubit reduction, verified element-wise in Kraus form, then
inducted across the QFT? A non-specialist-friendly channel-identity
appendix would make the result portable to circuit-QED / dynamic-circuit
engineers unfamiliar with consistent-histories notation.

\item[\textbf{Q5.}]\textbf{Deferred-measurement premium on nearest-neighbour connectivity.}
On hardware without \texttt{if\_test} support, the semiclassical variant
must be recompiled via the principle-of-deferred-measurement --- which
reintroduces two-qubit gates. Does the deferred-semiclassical circuit ever
yield a strictly shorter total depth than the original standard QFT after
transpilation onto a linear-nearest-neighbour coupling map at $n{=}10..20$?
If yes, quantify the dynamic-circuits premium; if no, formalize this as
evidence that the entire Griffiths-Niu benefit is contingent on
hardware-level classical feed-forward support.
\end{enumerate}
